%%%%%%%% ICML 2024 EXAMPLE LATEX SUBMISSION FILE %%%%%%%%%%%%%%%%%

\documentclass{article}

% Recommended, but optional, packages for figures and better typesetting:
\usepackage{microtype}
\usepackage{graphicx}
\usepackage{subfigure}
\usepackage{booktabs} % for professional tables

% hyperref makes hyperlinks in the resulting PDF.
% If your build breaks (sometimes temporarily if a hyperlink spans a page)
% please comment out the following usepackage line and replace
% \usepackage{icml2024} with \usepackage[nohyperref]{icml2024} above.
\usepackage{hyperref}


% Attempt to make hyperref and algorithmic work together better:
\newcommand{\theHalgorithm}{\arabic{algorithm}}

% Use the following line for the initial blind version submitted for review:
% \usepackage{icml2024}

% If accepted, instead use the following line for the camera-ready submission:
\usepackage[accepted]{icml2024}

% For theorems and such
\usepackage{amsmath}
\usepackage{amssymb}
\usepackage{mathtools}
\usepackage{amsthm}

% if you use cleveref..
\usepackage[capitalize,noabbrev]{cleveref}

%%%%%%%%%%%%%%%%%%%%%%%%%%%%%%%%
% THEOREMS
%%%%%%%%%%%%%%%%%%%%%%%%%%%%%%%%
\theoremstyle{plain}
\newtheorem{theorem}{Theorem}[section]
\newtheorem{proposition}[theorem]{Proposition}
\newtheorem{lemma}[theorem]{Lemma}
\newtheorem{corollary}[theorem]{Corollary}
\theoremstyle{definition}
\newtheorem{definition}[theorem]{Definition}
\newtheorem{assumption}[theorem]{Assumption}
\theoremstyle{remark}
\newtheorem{remark}[theorem]{Remark}

% Todonotes is useful during development; simply uncomment the next line
%    and comment out the line below the next line to turn off comments
%\usepackage[disable,textsize=tiny]{todonotes}
\usepackage[textsize=tiny]{todonotes}


%%%%%%%%%%%%%%%%%%%%%%%%%% Authors Customization %%%%%%%%%%%%%%%%%%%
\usepackage{algorithm}
\usepackage{algorithmic}
% \usepackage{caption}
\usepackage[bb=boondox]{mathalfa}
\usepackage{listings}
\usepackage{multirow,tabularx}
\usepackage{amssymb}
\usepackage[T1]{fontenc}
\usepackage{wrapfig}
\usepackage{subcaption}
\usepackage{multirow}
\usepackage[normalem]{ulem}
\AddToHook{begindocument/before}{\usepackage{lpsb-mcid}} % LPSB MCID - deferred load for hyperref compat
\useunder{\uline}{\ul}{}
%%%%%%%%%%%%%%%%%%%%%%%%%%%%%%%%%%%%%%%%%%%%%%%%%%%%%%%%%%%%%%%%%%%


%%%%%%%%%%%%%%%%%%%%%%%%%%%%%%%%%%%%%%%%%%%%%%%%%%%%%%%%%%%%%%%%%%%
\newcommand{\fix}{\marginpar{FIX}}
\newcommand{\new}{\marginpar{NEW}}
\newcommand{\framework}{OpenPAL} %{\textcolor{cyan}{[Undecided-Framework]}}
\newcommand{\Contra}{Contra}
\newcommand{\bfemph}[1]{\textbf{\textit{#1}}}
% \newcommand{\highlight}[1]{\textcolor{blue}{#1}}
\newcommand{\highlight}[1]{#1}
\newcommand{\ming}[1]{{\bf \color{orange} [[Ming says ``#1'']]}}
\newcommand{\qi}[1]{{\bf \color{blue} [[Qi says ``#1'']]}}
\newcommand{\zhai}[1]{{\bf \color{magenta} [[Zhai: ``#1'']]}}
\newcommand{\qis}[1]{{\bf \color{blue} [[Qi suggests ``#1'']]}}
\newcommand{\xian}[1]{{\bf \color{cyan} [[Fuxian says ``#1'']]}}
\newcommand{\chao}[1]{{\bf \color{red} [[Chao says ``#1'']]}}
%%%%%%%%%%%%%%%%%%%%%%%%%%%%%%%%%%%%%%%%%%%%%%%%%%%%%%%%%%%%%%%%%%%


% The \icmltitle you define below is probably too long as a header.
% Therefore, a short form for the running title is supplied here:
% \icmltitlerunning{Enhanced Adaptation from Environmental Context to Instruction in Open-Ended Embodied Agents}
\icmltitlerunning{Building Open-Ended Embodied Agent via Language-Policy Bidirectional Adaptation}

\begin{document}

\twocolumn[
% \icmltitle{Enhanced Adaptation from Environmental Context to Instruction in Open-Ended Embodied Agents}
\icmltitle{Building Open-Ended Embodied Agent via Language-Policy \\ Bidirectional Adaptation}

% It is OKAY to include author information, even for blind
% submissions: the style file will automatically remove it for you
% unless you've provided the [accepted] option to the icml2023
% package.

% List of affiliations: The first argument should be a (short)
% identifier you will use later to specify author affiliations
% Academic affiliations should list Department, University, City, Region, Country
% Industry affiliations should list Company, City, Region, Country

% You can specify symbols, otherwise they are numbered in order.
% Ideally, you should not use this facility. Affiliations will be numbered
% in order of appearance and this is the preferred way.
\icmlsetsymbol{equal}{*}

\begin{icmlauthorlist}
\icmlauthor{Shaopeng Zhai}{equal,shlab}
\icmlauthor{Jie Wang}{equal,shlab}
\icmlauthor{Tianyi Zhang}{equal,shlab}
\icmlauthor{Fuxian Huang}{equal,shlab}\\
\icmlauthor{Qi Zhang}{equal,shlab}
\icmlauthor{Ming Zhou}{equal,shlab}
\icmlauthor{Jing Hou}{equal,tju,shlab}
\icmlauthor{Yu Qiao}{shlab}
\icmlauthor{Yu Liu}{shlab}
%\icmlauthor{}{sch}
%\icmlauthor{}{sch}
\end{icmlauthorlist}

\icmlaffiliation{shlab}{Shanghai AI Laboratory}
\icmlaffiliation{tju}{Tongji University}
% \icmlaffiliation{sch}{School of ZZZ, Institute of WWW, Location, Country}

\icmlcorrespondingauthor{Shaopeng Zhai}{zhaishaopeng@pjlab.org.cn}
\icmlcorrespondingauthor{Yu Liu}{liuyu@pjlab.org.cn}

% You may provide any keywords that you
% find helpful for describing your paper; these are used to populate
% the "keywords" metadata in the PDF but will not be shown in the document
\icmlkeywords{Machine Learning, ICML}

\vskip 0.3in
]

% this must go after the closing bracket ] following \twocolumn[ ...

% This command actually creates the footnote in the first column
% listing the affiliations and the copyright notice.
% The command takes one argument, which is text to display at the start of the footnote.
% The \icmlEqualContribution command is standard text for equal contribution.
% Remove it (just {}) if you do not need this facility.

% \printAffiliationsAndNotice{}  % leave blank if no need to mention equal contribution
\printAffiliationsAndNotice{\icmlEqualContribution} % otherwise use the standard text.

% \doparttoc % Tell to minitoc to generate a toc for the parts
% \faketableofcontents % Run a fake tableofcontents command for the partocs

% \part{} % Start the document part
% \parttoc % Insert the document TOC

\input{sections/abstract}
\input{sections/intro}
\input{sections/preliminaries}
\input{sections/contra}
\input{sections/framework}
\input{sections/results}
% \input{sections/relatedwork}
\input{sections/conclusion}

% In the unusual situation where you want a paper to appear in the
% references without citing it in the main text, use \nocite

% \section*{Impact Statements}
% This paper presents work that is directed towards constructing open-ended embodied agents designed for applications in human-AI interaction, taking into account practical considerations, including high real-time requirements.
% There are many potential societal consequences of our work, none which we feel must be specifically highlighted here, as they are common in the AI community.
\section*{Author Contribution Statement}
The authors confirm their contribution as follows:

\textbf{Shaopeng Zhai}: team leadership, open-ended learning, LLM/RLAF training, agent analysis, architecture design\\
\textbf{Jie Wang}: infrastructure/framework engineering, non-goal agent training, open-ended learning, ablation studies, feature engineering\\
\textbf{Tianyi Zhang}: non-goal agent training, open-ended learning, feature engineering\\
\textbf{Fuxian Huang}: non-goal agent training, paper writing, open-ended learning\\
\textbf{Qi Zhang}: LLM training, RLAF training, paper writing, ablation studies\\
\textbf{Ming Zhou}: co-training framework, curriculum research, paper writing\\
\textbf{Jing Hou}: LLM training, paper writing
% \newpage
\bibliography{ref}
\bibliographystyle{icml2024}

\appendix
\onecolumn
\input{sections/appendix_one_in_all}

% \input{sections/app_secs/LLM_draft}

\end{document}


% This document was modified from the file originally made available by
% Pat Langley and Andrea Danyluk for ICML-2K. This version was created
% by Iain Murray in 2018, and modified by Alexandre Bouchard in
% 2019 and 2021 and by Csaba Szepesvari, Gang Niu and Sivan Sabato in 2022.
% Modified again in 2023 by Sivan Sabato and Jonathan Scarlett.
% Previous contributors include Dan Roy, Lise Getoor and Tobias
% Scheffer, which was slightly modified from the 2010 version by
% Thorsten Joachims & Johannes Fuernkranz, slightly modified from the
% 2009 version by Kiri Wagstaff and Sam Roweis's 2008 version, which is
% slightly modified from Prasad Tadepalli's 2007 version which is a
% lightly changed version of the previous year's version by Andrew
% Moore, which was in turn edited from those of Kristian Kersting and
% Codrina Lauth. Alex Smola contributed to the algorithmic style files.
